\documentclass{article}

\usepackage{arxiv}

\usepackage[utf8]{inputenc}
\usepackage[T1]{fontenc}
\usepackage{hyperref}
\usepackage{url}
\usepackage{booktabs}
\usepackage{amsfonts}
\usepackage{amssymb}
\usepackage{amsmath}
\usepackage{nicefrac}
\usepackage{microtype}
\usepackage{cleveref}
\usepackage{graphicx}
\usepackage{natbib}
\usepackage{doi}
\usepackage{tikz}

%% ============================================================
%% Reference parameters (derived values are computed by pgfmath)
%% ============================================================
% --- Inputs ---
\pgfmathsetmacro{\refWACC}{0.07}            % Weighted-average cost of capital
\pgfmathsetmacro{\refN}{30}                 % Plant operating lifetime [yr]
\pgfmathsetmacro{\refTc}{6}                 % Construction time [yr]
\pgfmathsetmacro{\refG}{0.02}               % Inflation rate
\pgfmathsetmacro{\refFindirect}{0.20}       % CAS30 indirect cost fraction
\pgfmathsetmacro{\refLandIntensity}{0.25}   % Land intensity [acres/MWe]
\pgfmathsetmacro{\refLandCost}{10000}       % Land cost [$/acre]
\pgfmathsetmacro{\refCbase}{100}            % Levelization example base cost [M$]
% --- Derived ---
\pgfmathsetmacro{\refCRF}{\refWACC*(1+\refWACC)^\refN / ((1+\refWACC)^\refN - 1)}
\pgfmathsetmacro{\refIDCfrac}{((1+\refWACC)^\refTc - 1)/(\refWACC*\refTc) - 1}
\pgfmathsetmacro{\refIDCpct}{\refIDCfrac * 100}
\pgfmathsetmacro{\refAone}{\refCbase * (1+\refG)^\refTc}
\pgfmathsetmacro{\refPV}{\refAone * (1 - ((1+\refG)/(1+\refWACC))^\refN) / (\refWACC - \refG)}
\pgfmathsetmacro{\refLevelized}{\refCRF * \refPV}
\pgfmathsetmacro{\refUnderestimate}{(\refLevelized - \refAone) / \refAone * 100}
\pgfmathsetmacro{\refLandTotal}{\refLandIntensity * \refLandCost / 1000}  % = acres/MWe * $/acre * 1000 MWe / 1e6

\title{A Differentiable Costing Framework for Fusion Power Plants}

\newif\ifuniqueAffiliation
\uniqueAffiliationtrue

\ifuniqueAffiliation
%\author{
%	Tal Rubin \\
%	Astera Institute \\
%	\texttt{tal.rubin@astera.org} \\
%	\And
%	Mallory Snowden \\
%	Astera Institute \\
%	\texttt{mallory.snowden@astera.org} \\
%	\And
%	Reid Westwood \\
%	Astera Institute \\
%	\texttt{westwood.reid@gmail.com} \\
%	\And
%	Damien Scott \\
%	Astera Institute \\
%	\texttt{damien.scott@astera.org} \\
%}
\fi

\renewcommand{\shorttitle}{Differentiable Fusion Costing}

\hypersetup{
pdftitle={A Differentiable Costing Framework for Fusion Power Plants},
pdfauthor={Tal Rubin, Mallory Snowden, Reid Westwood, Damien Scott},
pdfkeywords={fusion energy, LCOE, cost estimation, automatic differentiation, JAX},
}

\begin{document}
\maketitle

\begin{abstract}
We present \textsc{1costingfe}, a JAX-based costing framework for fusion power plants
that computes the levelized cost of electricity (LCOE) from customer-level requirements
across all major confinement families and fuel cycles.
The framework implements the Code of Accounts Structure (CAS)
established by the Generation~IV Economic Modeling Working Group and adapted for fusion 
by Schulte et~al.\ and the ARIES program, and adopted by the ARPA-E-funded pyfecons framework
of Woodruff Scientific.
The purpose of \textsc{1costingfe}, beyond enabling extraction of exact gradients and built-in uncertainty analysis via JAX,
is to assist the identification of technological corridors into truly competitive fusion power in the 1cent/kWh range.
We describe the economics module, the physics module, and the cost
account methodology, noting where our treatment diverges from prior
work.
\end{abstract}

\keywords{fusion energy \and LCOE \and cost estimation \and automatic differentiation \and JAX}


%% ============================================================
\section{Introduction}
\label{sec:intro}

Fusion energy promises abundant, low-carbon baseload power, but its economic
viability remains an open question. Credible cost estimates are essential for
guiding R\&D investment, comparing confinement concepts, and identifying the
design parameters with the greatest leverage on the levelized cost of
electricity (LCOE).

The most widely used costing methodology for fusion power plants is the
Code of Accounts Structure (CAS) developed by Schulte et~al.\
\citep{schulte1978} at Pacific Northwest Laboratory and later adopted by
the ARIES program \citep{miller1998} and the Generation~IV Economic Modeling
Working Group \citep{emwg2007}. This framework decomposes plant costs into
hierarchical accounts covering pre-construction, direct capital, indirect
services, financial charges, and annual operating costs. The PyFECONS
code \citep{woodruff2026, woodruff2026b} provides a Python implementation of this
methodology.

Power plant costing estimates are challenging even for mature technologies
with extensive historical data, with fission plants exhibiting significant
cost overruns and schedule delays \citep{eashgates2020}. Applying this
framework to fusion introduces two further challenges:
\begin{enumerate}
\item No combination of confinement approach (magnetic, inertial,
  magneto-inertial), system (tokamak, stellarator, laser-driven, etc.) 
  and fuel cycle (DT, DD, D${}^3$He, or p${}^{11}$B) has yet
  demonstrated sustained net electricity generation, and it remains unclear
  which will prove commercially viable.
\item Most proposed fusion plant designs are pre-conceptual, with
  incomplete physics models and unfinalized engineering geometry, making
  the subsystem and material inputs to the cost accounts inherently
  uncertain.
\end{enumerate}

To pursue this search systematically, we are developing \textsc{1costingfe}, a
costing framework built on JAX that computes LCOE from customer-level
requirements across all major confinement families and fuel cycles. The
choice of JAX as the computational backend provides automatic
differentiation - exact partial derivatives of LCOE with respect to every
input parameter in a single backward pass - and vectorized evaluation over
parameter sweeps and Monte Carlo samples. These capabilities transform
sensitivity analysis from a computational bottleneck into a routine
operation, enabling rapid identification of the design levers with the
greatest leverage on cost. 

We are entering a phase of fusion development where fusion power plants are being designed,
not only experiments meant to determine the plasma physics and demonstrate scientific breakeven.
A commercially viable fusion power plant must be designed with its economics in mind, which introduces
additional parameters and constraints that must be added to the optimization loop.
The economic parameters---construction time, economic plant lifetime, capital cost and operating costs,
and the cost of capital, to name a few---become critical components in the calculation, in addition to the
particular plasma physics, nuclear engineering, and traditional engineering parameters.

Most open-source fusion costing frameworks are either
\begin{enumerate}
\item physics-agnostic, accepting user inputs for many, if not all, cost
  accounts, or
\item applicable to a single, relatively mature confinement family and fuel
  cycle (e.g.\ D-T tokamaks), using reduced-order plasma models to
  determine fusion power from the coils and geometry and to size subsystems
  such as neutral beam injectors and blankets.
\end{enumerate}
Individual companies often maintain proprietary models that are unavailable
to the public, as they contain sensitive information about their design and
its maturity.

Certain inputs to the cost model are well constrained regardless of design
choice. The fuel nuclear physics: cross sections, energy per reaction, and 
products are immutable. Raw material isotope mix and enrichment ease are 
well characterized and unlikely to shift dramatically, though some cost 
reduction through learning is expected if a fusion plant design is successful.
Materials science for fusion-specific alloys is expected to lag behind deployment,
owing to the current lack of testing facilities for neutron irradiation and high heat fluxes,
but a learning curve is expected as demand for high-performance materials grows.

On the economics side, the cost of capital is a market condition rather
than a design parameter, and a high cost of capital strongly incentivizes
designs with lower capital costs and faster construction, even at the
expense of higher operating costs.

The \textsc{1cfe} initiative to search for corridors through which fusion energy can reach the 
optimistic price target of 1~cent/kWh must balance technological optimism regarding the 
potential of medium-to-high Technology Readiness Level (TRL) concepts with healthy skepticism
regarding very early stage physics. As an example, untested direct energy conversion methods are
not specifically considered at this stage, but the potential for their development and deployment is not ruled out.

Data availability for fusion power plant costs is low: no commercial plant
has been built, and the few large-scale projects (ITER, NIF) are research
facilities with non-representative cost structures. In this regime, a model
with too many free parameters risks overfitting to sparse calibration points
and misrepresenting the accuracy of its estimates. The framework is therefore
built bottom-up with deliberate parsimony: each cost account is independently
justified from literature and first principles, parameters are introduced
only when supported by data or physical reasoning, and methodology choices
are documented in per-account documentation.
The computation flows from customer requirements (net electric
power, availability, plant lifetime) through power balance, engineering
sizing, per-account costing, and financial aggregation to a single LCOE
figure, supporting cost overrides at any account level to accommodate
concept-specific engineering detail as it becomes available.

The framework is designed to integrate with other tools in the 1cfe ecosystem,
including \textsc{Fusion-TEA}, a SysML~v2 power plant model that can override
costing parameters for specific implementations, and a backcasting dashboard
for target-driven design exploration.

The remainder of this paper is organized as follows.
\Cref{sec:methodology} presents the economics module, the physics
module, and the cost account methodology, with detailed treatment of
the accounts for which we have completed independent justification
studies.


%% ============================================================
\section{Methodology}
\label{sec:methodology}

\subsection{Economics Module}
\label{sec:lcoe}

The levelized cost of electricity is the constant electricity price at
which the present value of revenue equals the present value of all costs
over the plant lifetime. It aggregates annualized capital and operating
costs into a single metric in \$/MWh:

\begin{equation}
\label{eq:lcoe}
\text{LCOE} = \frac{\text{Annualized cost}}{\text{Annual electricity production}}
\end{equation}

\noindent The annualized cost is the sum of annualized capital cost
$C_{\text{annual}}$ and annualized operating cost $M_{\text{annual}}$
(maintenance, fuel, and manpower), both in M\$/yr. The annual electricity
production is:

\begin{equation}
\label{eq:e-annual}
E_{\text{annual}} = 8760[\text{h}/\text{yr}] \times P_{\text{net}} [\text{MW}/\text{module}] \times N_{\text{mod}} \times A
\end{equation}

\noindent where $P_{\text{net}}$ is the net electric power per module (MW),
$N_{\text{mod}}$ is the number of reactor modules, $A$ is the plant
availability (capacity factor), and 8760 is the number of hours per year.
The computation of $P_{\text{net}}$ from fusion power and
engineering coefficients is described in the physics module
(\cref{sec:physics-module}).

\Cref{fig:cashflow} illustrates the cash-flow structure. Total capital
expenditure (CAPEX) is disbursed during the construction period, after
which revenue from electricity sales and annual operating expenditure
(OPEX) flow in opposite directions over the plant lifetime. 

\begin{figure}[htbp]
\centering
\begin{tikzpicture}[x=0.52cm, y=0.45cm]
  % --- timeline ---
  \draw[thick] (0,0) -- (22,0);

  % --- construction and licensing period bracket ---
  \draw[thick, stealth-stealth, color=blue!70!black]
    (0,-5.5) -- node[midway, below, font=\small] {Construction time $T_c$} (8,-5.5);

  % --- plant lifetime bracket ---
  \draw[thick, stealth-stealth, color=blue!70!black]
    (8,-5.5) -- node[midway, below, font=\small] {Plant lifetime $n$ years} (22,-5.5);

  % --- CAPEX arrows (uniform disbursements during construction) ---
  \foreach \x in {0,2,4,6} {
    \draw[thick, -{stealth}, color=black!80] (\x,0) -- (\x,-3.5);
  }
  \node[font=\small] at (3,-4.2) {CAPEX};

  % --- Revenue arrows (upward, during operation, spacing=2, starting 2 from COD) ---
  \foreach \x in {10,12,14,16,18,20,22} {
    \draw[thick, -{stealth}, color=green!50!black] (\x,0) -- (\x,3.5);
  }
  \node[font=\small, color=green!50!black] at (16,4.2) {Revenue};

  % --- OPEX arrows (downward, growing to reflect nominal escalation) ---
  \foreach \x/\h in {10/-1.4, 12/-1.6, 14/-1.8, 16/-2.0, 18/-2.2, 20/-2.5, 22/-2.8} {
    \draw[thick, -{stealth}, color=black!80] (\x,0) -- (\x,\h);
  }
  \node[font=\small] at (16,-3.4) {OPEX};

  % --- time markers ---
  \draw[thin, gray] (8,-0.3) -- (8,0.3);
  \node[font=\scriptsize, gray, above] at (8,0.15) {COD};
\end{tikzpicture}
\caption{Schematic cash-flow diagram for a fusion power plant.
  Capital expenditure (CAPEX) is disbursed uniformly during the
  construction period $T_c$, accruing interest until the commercial
  operation date (COD). After COD, revenue from electricity sales and
  annual operating expenditure (OPEX) flow over the plant operating
  lifetime of $n$ years. OPEX grows in nominal terms at the inflation
  rate. The LCOE (\cref{eq:lcoe}) is the constant electricity price at
  which the present value of revenue equals the present value of all
  costs.}
\label{fig:cashflow}
\end{figure}

\subsubsection{Capital Recovery Factor}

The Capital Recovery Factor (CRF) converts a present-value capital
investment into equal annual payments over the plant lifetime:

\begin{equation}
\label{eq:crf}
\text{CRF}(i, n) = \frac{i\,(1+i)^n}{(1+i)^n - 1}
\end{equation}

\noindent where $i$ is the weighted-average cost of capital (WACC) and
$n$ is the plant operating lifetime in years. At reference conditions
($i = \refWACC$, $n = \pgfmathprintnumber[precision=0]{\refN}$~yr),
CRF $= \pgfmathprintnumber[precision=4]{\refCRF}$.

\subsubsection{Annualized Capital Cost}
\label{sec:capex-annual}

Capital expenditure is disbursed over the construction period $T_c$
before any revenue is generated (\cref{fig:cashflow}). Assuming the
overnight capital cost $C$ is spent in $T_c$ equal annual installments
of $C/T_c$, each installment accrues compound interest from the time
it is spent until the commercial operation date (COD). The future value
of this series at COD is:

\begin{equation}
\label{eq:fv-capex}
\text{FV} = \frac{C}{T_c} \sum_{k=0}^{T_c - 1} (1+i)^{k}
          = \frac{C}{T_c} \cdot \frac{(1+i)^{T_c} - 1}{i}
\end{equation}

\noindent The interest during construction (IDC) is the excess above
the overnight cost---the total financing charge incurred while the
plant earns no revenue:

\begin{equation}
\label{eq:idc-lcoe}
\text{IDC} = \text{FV} - C
           = C \left[ \frac{(1+i)^{T_c} - 1}{i\,T_c} - 1 \right]
\end{equation}

\noindent The total capital investment at COD is therefore
$C + \text{IDC} = \text{FV}$, and the annualized capital cost is:

\begin{equation}
\label{eq:c-annual}
C_{\text{annual}} = \text{CRF}(i, n) \times (C + \text{IDC})
\end{equation}

\noindent At reference conditions ($i = \refWACC$,
$T_c = \pgfmathprintnumber[precision=0]{\refTc}$~yr), the IDC adds
\pgfmathprintnumber[precision=1]{\refIDCpct}\% to the overnight cost.

\subsubsection{Levelized Annual Operating Cost}
\label{sec:lac}

Annual operating costs (maintenance, fuel, and manpower) are stated in constant
base-year dollars but escalate at the inflation rate $g$ in nominal
terms over the plant lifetime. The levelized annual cost accounts for
this growth using a three-step procedure.

First, the base-year annual cost $C_0$ is inflated to
first-year-of-operation dollars:

\begin{equation}
\label{eq:inflate}
A_1 = C_0 \, (1 + g)^{T_c}
\end{equation}

\noindent where $T_c$ is the construction time (or total project time
for first-of-a-kind plants, which includes licensing).

Second, the present value of the resulting growing annuity is:

\begin{equation}
\label{eq:growing-annuity}
\text{PV} =
\begin{cases}
\displaystyle A_1 \cdot \frac{1 - \left(\frac{1+g}{1+i}\right)^n}{i - g}
  & \text{if } i \neq g \\[12pt]
\displaystyle A_1 \cdot \frac{n}{1+i}
  & \text{if } i = g
\end{cases}
\end{equation}

Third, the levelized annual cost is obtained by applying the CRF:

\begin{equation}
\label{eq:levelized}
C_{\text{level}} = \text{CRF}(i, n) \times \text{PV}
\end{equation}

\noindent The $i = g$ case follows from L'H\^{o}pital's rule and is
implemented using \texttt{jnp.where} for JAX traceability.
For example, at $C_0 = \pgfmathprintnumber[precision=0]{\refCbase}$~M\$,
$i = \refWACC$, $g = \refG$,
$n = \pgfmathprintnumber[precision=0]{\refN}$~yr,
$T_c = \pgfmathprintnumber[precision=0]{\refTc}$~yr, the levelized
cost is \pgfmathprintnumber[precision=1]{\refLevelized}~M\$/yr
compared to \pgfmathprintnumber[precision=1]{\refAone}~M\$/yr from the
na\"{\i}ve single-year inflation formula---a
\pgfmathprintnumber[precision=0]{\refUnderestimate}\% underestimate
when the growing annuity is neglected.


%% ------------------------------------------------------------
\subsection{Physics Module}
\label{sec:physics-module}

The denominator of the LCOE (\cref{eq:lcoe}) is the annual saleable
electricity, which depends on the net electric power $P_{\text{net}}$ and plant availability.
The purpose of the physics module is computing $P_{\text{net}}$ from
device and plasma parameters, given a confinement concept and fuel cycle.
A computation of the plant availability would also fit in this, or an Engineering module, 
but requires a dedicated model incorporating the specifics of a design (e.g. maintainability, accessibility, component failure modes).

Currently, the physics module is under development, and has some sections
already implemented. The power balance section is complete and concept-agnostic,
while the plasma performance section is being developed in parallel with the 
first confinement-specific layer: a zero-dimensional tokamak equilibrium model.

The power balance section of the physics module computes $P_{\text{net}}$ 
from the fusion power $P_{\text{fus}}$ for the four fusion fuels by tracking
every energy channel from the plasma to the grid: energy split between fusion ash and neutrons, 
radiation losses, blanket multiplication, thermal and direct energy conversion, and
recirculating power. 

Many concepts will deviate significantly from the defaults - a non-Maxwellian 
deuterium plasma may alter the burn rates of the secondary D-D reactions,
or a doped plasma may alter the radiation losses or the plasma optical thickness
\citep{frank2022}. So these and other coefficients are exposed as first-class 
parameters amenable to sensitivity analysis via automatic differentiation.

Lumped efficiency parameters are used for characterizing the power flow and thermal cycle.
A more involved power plant thermodynamics model could be implemented in the future,
which would allow the selection of a thermal cycle (Bryton vs. Rankine, etc.) and working 
fluid. 

The power balance above is concept-agnostic: given $P_{\text{fus}}$
and the engineering coefficients, it produces $P_{\text{net}}$
regardless of confinement scheme. To close the loop and determine
$P_{\text{fus}}$ from reactor geometry and plasma parameters, a
concept-specific physics model is required.
\Cref{sec:tokamak-model} describes the zero-dimensional tokamak
equilibrium model included as the first such layer; models for other
confinement families will follow the same pattern.

The power balance described here applies to magnetic fusion energy
(MFE) systems; inertial and magneto-inertial variants follow the same
logic with minor differences noted below.

\subsubsection{Fusion Power Split}

The fusion power divides into charged products (ash) and neutrons
according to the reaction kinematics of the fuel cycle. For D-T
fusion, each reaction releases 17.6~MeV: 3.5~MeV to the alpha
particle and 14.1~MeV to the neutron, giving a neutron fraction
of approximately 0.80. Aneutronic fuels such as p-${}^{11}$B
produce only charged particles (neutron fraction~$\approx 0$).
\Cref{tab:neutron-fractions} summarizes the primary split for each
fuel cycle.

\begin{table}[htbp]
\caption{Fusion power split by fuel cycle (primary reactions only;
  secondary burn of tritium and ${}^3$He products is included in the
  full calculation).}
\label{tab:neutron-fractions}
\centering
\begin{tabular}{lccl}
\toprule
Fuel & $Q$ (MeV) & $f_n$ & Products \\
\midrule
D-T           & 17.6  & 0.80 & ${}^4$He + n \\
D-D           & 3.65  & 0.33 & ${}^3$He + n \\
              &       &      & T + p \\
D-${}^3$He    & 18.3  & 0.0  & ${}^4$He + p \\
p-${}^{11}$B  & 8.7   & 0.0  & 3${}^4$He \\
\bottomrule
\end{tabular}
\end{table}

\paragraph{D-D secondary burn.}
The D-D fuel cycle is more complex than the table suggests because both
branches produce fusion fuels as ash. Branch~1 produces tritium and
branch~2 produces ${}^3$He, each of which can undergo secondary fusion
with the background deuterium:

\begin{align}
\label{eq:dd-secondary}
\text{T} + \text{D} &\to {}^4\text{He}\,(3.5\;\text{MeV})
  + \text{n}\,(14.1\;\text{MeV}), &\quad Q_{\text{DT}} &= 17.6\;\text{MeV} \\
{}^3\text{He} + \text{D} &\to {}^4\text{He}\,(3.6\;\text{MeV})
  + \text{p}\,(14.7\;\text{MeV}), &\quad Q_{\text{D}{}^3\text{He}} &= 18.3\;\text{MeV}
\end{align}

\noindent The fraction of ash that undergoes secondary burn depends on
the plasma temperature and confinement time; the framework parameterizes
this through burn fractions $f_T$ and $f_{{}^3\text{He}}$. The effective
energy per D-D event, averaged over the two primary branches, becomes:

\begin{equation}
\label{eq:dd-etotal}
\bar{E}_{\text{total}} = \tfrac{1}{2}(Q_1 + Q_2)
  + \tfrac{1}{2} f_T \, Q_{\text{DT}}
  + \tfrac{1}{2} f_{{}^3\text{He}} \, Q_{\text{D}{}^3\text{He}}
\end{equation}

\noindent where $Q_1 = 4.03$~MeV and $Q_2 = 3.27$~MeV are the primary
branch Q-values. The effective charged-particle energy follows the same
structure, noting that the secondary D-${}^3$He reaction produces only
charged particles while the secondary D-T reaction splits between an
alpha and a neutron. At default burn fractions ($f_T = 0.97$,
$f_{{}^3\text{He}} = 0.69$), the effective total energy rises from
3.65~MeV to approximately 18.5~MeV per primary event, and the effective
neutron fraction increases from 0.33 to approximately 0.44 due to the
secondary D-T neutrons.

\paragraph{D-${}^3$He side reactions.}
The D-${}^3$He fuel cycle is nominally aneutronic, but at the
temperatures required for D-${}^3$He fusion, a fraction $f_{\text{DD}}$
of the deuterium undergoes D-D side reactions. These side reactions
produce tritium (branch~1) which can in turn undergo secondary D-T
fusion, introducing 14.1~MeV neutrons into an otherwise neutron-free
system. The effective ash fraction is the weighted average of the pure
D-${}^3$He contribution (fully charged) and the D-D side-reaction
contribution (with its own secondary burn chain):

\begin{equation}
\label{eq:dhe3-ash}
f_{\text{ash}} = (1 - f_{\text{DD}})
  + f_{\text{DD}} \cdot \frac{\bar{E}_{\text{charged,DD}}}
                              {\bar{E}_{\text{total,DD}}}
\end{equation}

\noindent where $\bar{E}_{\text{charged,DD}}$ and
$\bar{E}_{\text{total,DD}}$ include the secondary D-T burn of the
tritium product (the ${}^3$He product from the D-D side reaction is
already present in excess and need not be tracked separately). At
the default D-D fraction $f_{\text{DD}} = 0.07$ and tritium burn
fraction $f_T = 0.97$, the effective neutron fraction is approximately
3\%, small but nonzero---and consequential for shielding, activation,
and tritium breeding requirements.

\paragraph{p-${}^{11}$B side reactions.}
The p-${}^{11}$B reaction is nominally aneutronic, but two side channels
produce neutrons. The ${\sim}2.9$~MeV alpha particles from the primary
reaction can undergo ${}^{11}$B($\alpha$,n)${}^{14}$N
($Q = +0.16$~MeV), which is exothermic but suppressed by the Coulomb
barrier (${\sim}3.2$~MeV); measured cross sections at alpha energies of
2.5--3.0~MeV are 200--600~mb with significant resonance structure
\citep{walker1949, prior2017}. The endothermic channel
${}^{11}$B(p,n)${}^{11}$C ($Q = -2.76$~MeV, threshold 3.02~MeV) is
accessible to protons in the high-energy tail, though at a much lower
rate (${\sim}10^{-5}$ relative to the primary reaction). Published
estimates place the overall neutron fraction at approximately 0.1\% of
reactions, carrying less than 0.2\% of the total released energy
\citep{putvinski2019}.

The framework models these side reactions via burn fractions
$f_{\alpha n}$ and $f_{pn}$, analogous to the secondary-burn treatment
used for D-D fuel. Both default to zero: \citet{ochs2025} show that
viable p-${}^{11}$B reactors must remove alpha ash on timescales much
shorter than the energy confinement time to avoid bremsstrahlung
poisoning, which drastically reduces the population of MeV-scale alphas
coexisting with the boron fuel and thus suppresses the dominant
${}^{11}$B($\alpha$,n)${}^{14}$N channel
\citep{hora2021}.


\subsubsection{Thermal and Electric Power}

The charged-particle (ash) power divides into radiation losses
(bremsstrahlung and synchrotron) and transport power deposited on
plasma-facing surfaces. The framework computes radiation power from
plasma parameters assuming an optically thin plasma
\citep{wesson2011}:

\begin{equation}
\label{eq:p-brem}
P_{\text{brem}} = 5.35 \times 10^{-37}\; n_e^2 \, Z_{\text{eff}}
  \sqrt{T_e}\; V \quad [\text{W}]
\end{equation}

\begin{equation}
\label{eq:p-sync}
P_{\text{sync}} = 6.2 \times 10^{-22}\; n_e \, T_e^2 \, B^2 \; V
  \quad [\text{W}]
\end{equation}

\noindent where $n_e$ is the electron density in m$^{-3}$, $T_e$ the
electron temperature in keV, $Z_{\text{eff}}$ the effective charge,
$B$ the magnetic field in T, and $V$ the plasma volume in m$^3$.
Synchrotron radiation applies only to MFE; IFE and MIF set
$P_{\text{rad}} = 0$. This is a naive model: optically thick plasmas
or those with significant wall reflection would re-absorb a fraction
of the synchrotron radiation, reducing the effective loss.

\paragraph{Impurity line radiation.}
In MFE plasmas, line radiation from partially ionized impurities is
often the dominant radiation loss channel. The framework models this
via a steady-state sputtering--transport chain connecting the
first-wall material to the core impurity concentration and radiative
loss.

The wall material is specified as one of six options: tungsten~(W),
carbon~(C), beryllium~(Be), molybdenum~(Mo), silicon carbide~(SiC),
or liquid lithium~(Li). For each material, a simplified Bohdansky
physical sputtering yield \citep{bohdansky1984, eckstein2007} gives
the number of wall atoms released per incident ion at the
sheath-accelerated energy $E_{\text{ion}} \approx 3\,T_{\text{edge}}$:

\begin{equation}
\label{eq:sputtering}
Y(E) = Q \left(1 - \left(\frac{E_{\text{th}}}{E}\right)^{2/3}\right)
       \left(1 - \frac{E_{\text{th}}}{E}\right)^{2}
\end{equation}

\noindent where $E_{\text{th}}$ is the material-dependent threshold
energy and $Q$ is a fit coefficient. Below threshold ($E < E_{\text{th}}$),
$Y = 0$. The steady-state impurity fraction in the core is then:

\begin{equation}
\label{eq:f-impurity}
f_z = Y(3\,T_{\text{edge}}) \cdot
  \frac{A_{\text{fw}}}{V} \cdot
  \frac{\tau_{\text{imp}}}{\tau_E} \cdot f_{\text{screen}}
\end{equation}

\noindent where $A_{\text{fw}}/V$ is the first-wall area to plasma
volume ratio, $\tau_{\text{imp}}/\tau_E$ is the impurity-to-energy
confinement time ratio (default 3.0), and $f_{\text{screen}}$ is the
scrape-off layer screening factor (default 0.01 for high-$Z$
materials, 0.1 for low-$Z$). Lithium is a special case: as a fully
ionized species above ${\sim}0.1$~keV, it produces negligible line
radiation at reactor core temperatures; the primary benefit of a
liquid Li wall is low sputtering yield from continuous surface renewal
and correspondingly low $f_{\text{screen}}$.

Deliberately seeded impurities (e.g.\ neon, argon) for divertor
radiation or edge cooling can be specified directly via their
concentration $f_z$, bypassing the sputtering model.

The impurity line radiation power uses coronal-equilibrium radiative
loss rate coefficients $L_z(T_e)$
\citep{post1977, mavrin2018}, represented as piecewise
power-law fits for each species:

\begin{equation}
\label{eq:p-line}
P_{\text{line}} = n_e^2 \sum_z f_z \, L_z(T_e) \; V \quad [\text{W}]
\end{equation}

\noindent The summation runs over all impurity species (both
wall-derived and seeded). Cooling curves are tabulated for W, C, Be,
Mo, Si, Li, Ne, and Ar. The total
radiated power is capped at the ash power,
$P_{\text{rad}} = \min(P_{\text{brem}} + P_{\text{sync}} + P_{\text{line}},\; P_{\text{ash}})$,
and the remainder is the transport power
$P_{\text{wall}} = P_{\text{ash}} - P_{\text{rad}}$.

All new parameters default to no-op values (null wall material, empty
seeded impurities), preserving backward compatibility: when no wall
material is specified, $P_{\text{line}} = 0$ and the radiation model
reduces to bremsstrahlung plus synchrotron as before.

Neutron power is absorbed in the breeding blanket, where nuclear
reactions (primarily ${}^6$Li(n,$\alpha$)T) release additional energy
characterized by the blanket multiplication factor $M_n$ (typically
1.1 for lithium blankets). The total thermal power delivered to the
heat-transfer system is:

\begin{equation}
\label{eq:p-thermal}
P_{\text{th}} = M_n \, P_n + P_{\text{rad}} + P_{\text{wall}}
              + P_{\text{aux}} + \eta_p \, P_{\text{pump}}
\end{equation}

\noindent where $P_n$ is the neutron power, $P_{\text{rad}}$ is the
radiated power reabsorbed by the first wall, $P_{\text{wall}}$ is the
charged-particle transport power deposited on plasma-facing components,
$P_{\text{aux}}$ is the injected auxiliary heating power, and
$\eta_p P_{\text{pump}}$ is the fraction of pumping power deposited as
heat. A conventional Rankine cycle with thermal efficiency $\eta_{\text{th}}$
converts thermal power to gross electric power:

\begin{equation}
\label{eq:p-gross}
P_{\text{gross}} = \eta_{\text{th}} \, P_{\text{th}}
\end{equation}

\noindent For systems with direct energy conversion (DEC), a fraction
$f_{\text{dec}}$ of the charged-particle transport power bypasses the
thermal cycle and is converted at efficiency $\eta_{\text{de}}$, and
the gross electric power becomes
$P_{\text{gross}} = \eta_{\text{th}} P_{\text{th}} +
 f_{\text{dec}} \, \eta_{\text{de}} \, P_{\text{transport}}$.
DEC is disabled by default ($f_{\text{dec}} = 0$) pending further
technology maturation.

\subsubsection{Recirculating Power and Net Electric Output}

A fraction of the gross electric power is recirculated to sustain
the plant:

\begin{equation}
\label{eq:p-recirc}
P_{\text{recirc}} = \frac{P_{\text{aux}}}{\eta_{\text{aux}}}
  + P_{\text{coils}} + P_{\text{pump}} + P_{\text{cryo}}
  + P_{\text{cool}} + f_{\text{sub}} \, P_{\text{gross}}
  + P_{\text{other}}
\end{equation}

\noindent where $P_{\text{aux}}/\eta_{\text{aux}}$ is the wall-plug
power for the heating system, $P_{\text{coils}}$ powers the magnets,
$P_{\text{cryo}}$ and $P_{\text{cool}}$ serve the cryogenic and cooling
systems, $f_{\text{sub}} P_{\text{gross}}$ covers balance-of-plant
subsystems (typically 3\%), and $P_{\text{other}}$ includes tritium
processing and housekeeping loads. The net electric power is:

\begin{equation}
\label{eq:p-net}
P_{\text{net}} = P_{\text{gross}} - P_{\text{recirc}}
\end{equation}

\noindent The engineering gain $Q_{\text{eng}} = P_{\text{gross}} /
P_{\text{recirc}}$ must exceed unity for the plant to produce net
electricity. The recirculating fraction
$1/Q_{\text{eng}}$ is a key figure of merit: lower recirculating
fractions leave more power for sale and directly reduce the LCOE
denominator.

For inertial fusion energy (IFE) systems, $P_{\text{aux}}$ is replaced
by driver power (laser or pulsed power) with its own wall-plug
efficiency, and target fabrication adds a fixed power load. For
magneto-inertial fusion (MIF), both driver and target/liner fabrication
power contribute. The power balance structure is otherwise identical.


%% ------------------------------------------------------------
\subsection{0D Tokamak Model}
\label{sec:tokamak-model}

The power balance of the physics module accepts fusion power
$P_{\text{fus}}$ as an input. To derive $P_{\text{fus}}$ from machine
geometry and plasma parameters for a tokamak, we implement a
zero-dimensional (0D) equilibrium model that chains standard scaling
laws. The model is fully JAX-differentiable, enabling gradient-based
sensitivity analysis of LCOE with respect to machine-level design
choices ($R$, $a$, $B$, $q_{95}$, etc.).

\subsubsection{DT Reactivity}

The DT fusion reactivity $\langle\sigma v\rangle$ is evaluated using
the analytic parameterization of \citet{boschhale1992}, valid for
$0.2 \leq T \leq 100$~keV:

\begin{equation}
\label{eq:sigma-v}
\langle\sigma v\rangle = C_1\,\theta\,
  \sqrt{\frac{\xi}{m_r c^2\,T^3}}\;\exp(-3\xi)
  \quad [\text{m}^3/\text{s}]
\end{equation}

\noindent where $\theta = T / (1 - T(C_2 + T(C_4 + TC_6))/(1 + T(C_3 + T(C_5 + TC_7))))$,
$\xi = (B_G^2 / 4\theta)^{1/3}$, $B_G = 34.38$~keV$^{1/2}$ is the
Gamow constant, and $C_1$--$C_7$ are the tabulated coefficients from
\citet{boschhale1992}. The reactivity peaks near $T \approx 65$~keV.

\subsubsection{Plasma Equilibrium}

Given major radius $R$, minor radius $a$, elongation $\kappa$,
toroidal field $B$, and edge safety factor $q_{95}$, the plasma
current follows from MHD force balance:

\begin{equation}
\label{eq:ip}
I_p = \frac{2\pi\,a^2\,\kappa\,B}{\mu_0\,R\,q_{95}}
  \quad [\text{A}]
\end{equation}

The electron density is set as a fraction $f_{\text{GW}}$ of the
Greenwald density limit \citep{iterphysics1999}:

\begin{equation}
\label{eq:greenwald}
n_{\text{GW}} = \frac{I_p\,[\text{MA}]}{\pi\,a^2}
  \quad [10^{20}\;\text{m}^{-3}],
\qquad
n_e = f_{\text{GW}}\,n_{\text{GW}}
\end{equation}

The plasma volume and first-wall area of the elongated torus are
$V = 2\pi^2 R\,a^2\,\kappa$ and $A_{\text{fw}} = 4\pi^2 R\,a\,\kappa$.

\subsubsection{Fusion Power}

For a 50/50 D-T plasma ($n_D = n_T = n_e/2$), the volumetric fusion
power is:

\begin{equation}
\label{eq:pfus-tokamak}
P_{\text{fus}} = \tfrac{1}{4}\,n_e^2\,
  \langle\sigma v\rangle(T_e)\;E_{\text{fus}}\;V
  \quad [\text{W}]
\end{equation}

\noindent where $E_{\text{fus}} = 17.58$~MeV per DT reaction.
This $P_{\text{fus}}$ feeds directly into the fuel-cycle-dependent
power split (\cref{tab:neutron-fractions}) and the power balance
of \cref{eq:p-thermal}--\cref{eq:p-net}.

\subsubsection{Energy Confinement}

The energy confinement time is evaluated against the IPB98(y,2)
H-mode scaling \citep{iterphysics1999}:

\begin{equation}
\label{eq:ipb98y2}
\tau_{E,98} = 0.0562\;
  I_p^{0.93}\,B^{0.15}\,\bar{n}_{19}^{0.41}\,
  P_{\text{heat}}^{-0.69}\,R^{1.97}\,\epsilon^{0.58}\,
  \kappa^{0.78}\,M^{0.19}
  \quad [\text{s}]
\end{equation}

\noindent where $\bar{n}_{19}$ is the line-averaged density in units of
$10^{19}$~m$^{-3}$, $P_{\text{heat}} = P_\alpha + P_{\text{aux}}$ is
the total heating power in MW, $\epsilon = a/R$ is the inverse aspect
ratio, and $M$ is the effective ion mass in AMU. The ratio of the
actual confinement time $\tau_E = W_{\text{th}} / P_{\text{heat}}$
(where $W_{\text{th}} = 3\,n_e\,T_e\,V$ for electrons and ions at
equal temperatures) to the scaling prediction gives the H-factor,
$H = \tau_E / \tau_{E,98}$, which measures the quality of confinement
relative to the empirical database.

\subsubsection{Stability and Engineering Limits}

The normalized beta,

\begin{equation}
\label{eq:beta-n}
\beta_N = \frac{\beta_t\,a\,B}{I_p\,[\text{MA}]},
\qquad
\beta_t = \frac{2\mu_0\,n_e\,T_e}{B^2},
\end{equation}

\noindent must remain below the Troyon limit ($\beta_N \lesssim 3.5$)
for ideal MHD stability. The safety factor $q_{95}$ must exceed
${\sim}2$ to avoid kink instabilities.

Two engineering constraints are evaluated: the neutron wall loading
$q_{\text{wall}} = P_n / A_{\text{fw}}$, which drives structural
damage and blanket lifetime; and the divertor heat flux, estimated
from a simplified scrape-off-layer model with midplane power width
$\lambda_q \approx 2$~mm and flux expansion factor $f_{\text{exp}}
\approx 5$.

\subsubsection{Disruption Model}

Disruptions destroy confinement and can cause significant structural
damage. The model estimates a disruption frequency from the proximity
of the operating point to stability boundaries:

\begin{equation}
\label{eq:disruption-rate}
\Gamma = \sum_{j} \Gamma_0\,\exp\!\left(-s\,m_j\right)
\end{equation}

\noindent where the sum runs over the Greenwald ($m = 1 - f_{\text{GW}}$),
Troyon ($m = 1 - \beta_N/\beta_{N,\max}$), and kink
($m = 1 - q_{95,\min}/q_{95}$) channels, $\Gamma_0 = 0.1$~FPY$^{-1}$
is the baseline rate far from limits, and $s = 15$ sets the
steepness. Each disruption damages a fraction of the first-wall
component lifetime and incurs 72~hours of downtime, reducing
effective plant availability.

\subsubsection{Operating Modes}

The model supports two operating modes. In \emph{forward mode}, the
user specifies machine geometry, field, safety factor, Greenwald
fraction, and temperature; the model returns the full plasma state
including $P_{\text{fus}}$, $\beta_N$, $H$-factor, wall loading,
and disruption rate. In \emph{inverse mode}, the user specifies a
target $P_{\text{net}}$ and the model solves for the electron
temperature $T_e$ that produces the required fusion power via
bisection, returning both the plasma state and the complete power
table. The inverse mode enables target-driven design: given a net
power requirement, the framework determines the plasma conditions,
checks them against stability and engineering limits, and computes
the resulting LCOE in a single differentiable pass.


%% ------------------------------------------------------------
\subsection{Cost Account Structure}
\label{sec:cas}

The cost account structure maps the financial framework of
\cref{sec:lcoe} onto specific plant systems and activities. Capital
accounts CAS10--CAS60 sum to the total capital investment (CAPEX);
annualized operating accounts CAS70 and CAS80 constitute OPEX; and
CAS90 annualizes the CAPEX via the CRF. \Cref{tab:cas-overview}
shows the full hierarchy.

\begin{table}[htbp]
\caption{Cost account structure overview. Accounts with dedicated
methodology subsections are marked with~$\dagger$.}
\label{tab:cas-overview}
\centering
\begin{tabular}{llp{7.0cm}}
\toprule
CAS & Description & Method \\
\midrule
10$^\dagger$ & Pre-construction & Land intensity scaling, fuel-dependent licensing \\
20 & Total direct costs & Sum of CAS21--29 \\
21 & Buildings \& structures & Power scaling on gross electric \\
22 & Reactor plant equipment & Component-level parametric (18 sub-accounts) \\
\hspace{1em}22.01.10$^\dagger$ & Remote handling \& maintenance & Fuel- and concept-dependent \\
\hspace{1em}22.01.12$^\dagger$ & Isotope separation plant & Zeroed; market purchase in CAS80 \\
23 & Turbine plant equipment & Power scaling on gross electric \\
24 & Electric plant equipment & Power scaling on gross electric \\
25 & Miscellaneous plant equip.\ & Power scaling \\
26 & Heat rejection systems & Thermal power scaling \\
27 & Special materials & Placeholder \\
28 & Digital twin & Fixed cost \\
29 & Contingency on direct costs & 10\% FOAK, 0\% NOAK \\
30$^\dagger$ & Indirect service costs & Fraction of CAS20, construction-time scaled \\
40 & Owner's costs & 5\% of CAS20 \\
50 & Supplementary costs & Spares, shipping, taxes, decommissioning \\
60$^\dagger$ & Interest during construction & Uniform-spending IDC formula \\
70$^\dagger$ & Annualized O\&M + replacement & Growing-annuity levelization \\
80 & Annualized fuel cost & Fuel-specific consumables, levelized \\
90$^\dagger$ & Annualized financial costs & CRF $\times$ total capital \\
\bottomrule
\end{tabular}
\end{table}

The following subsections detail the methodology for each account
that has undergone independent review and justification. Accounts
without a dedicated subsection use parametric scaling relations from
the ARIES/PyFECONS tradition; their documentation is ongoing.


%% ------------------------------------------------------------
\subsubsection{CAS10: Pre-Construction Costs}
\label{sec:cas10}

CAS10 covers land acquisition, site and plant permits, licensing,
engineering studies, and pre-construction reports. It is typically less
than 1\% of total capital; individual sub-accounts do not warrant
detailed parametric modeling.

\paragraph{Land and land rights (C11).}
Modern compact fusion plant designs require 20--200 acres, far less
than the legacy 1{,}000-acre assumption in earlier ARIES studies.
The CFS ARC project \citep{cfs_arc} plans 100 acres for 400~MWe,
yielding a land intensity of \refLandIntensity~acres/MWe. At US
industrial-zoned land costs of approximately
\$\pgfmathprintnumber[1000 sep={{\,}}, precision=0]{\refLandCost}/acre
(consistent with USDA 2025 farmland averages scaled for
industrial zoning), land cost for a 1~GWe plant is approximately
\$\pgfmathprintnumber[precision=1]{\refLandTotal}M---negligible in the context of multi-billion-dollar capital
investments.

The compact siting is enabled by the absence of an exclusion zone:
the NRC's 2023 decision to regulate D-T fusion under 10~CFR~Part~30
(byproduct materials) rather than Part~50/52 (reactor licensing)
eliminates the large-radius emergency planning zone required for
fission plants \citep{nrc2023}. Aneutronic fuels such as p-${}^{11}$B
likely fall outside NRC jurisdiction entirely.

\paragraph{Plant licensing (C13).}
Licensing timelines are fuel-dependent and affect first-of-a-kind
(FOAK) builds by extending the total project time, which flows into
the levelized annual cost calculation for CAS70 and CAS80.
\Cref{tab:licensing} summarizes the adopted values.

\begin{table}[htbp]
\caption{Licensing timeline by fuel cycle.}
\label{tab:licensing}
\centering
\begin{tabular}{lll}
\toprule
Fuel & Licensing time (yr) & Rationale \\
\midrule
D-T   & 2.0 & Part~30 licensing, upper end of 1--2~yr range \\
D-D   & 1.5 & Reduced tritium inventory \\
D-${}^3$He & 0.75 & Minimal radioactivity \\
p-${}^{11}$B & 0.0 & No NRC jurisdiction \\
\bottomrule
\end{tabular}
\end{table}


%% ------------------------------------------------------------
\subsubsection{CAS22.01.10: Remote Handling \& Maintenance Equipment}
\label{sec:cas2210}

CAS22.01.10 covers the robotic and remotely-operated equipment required
for in-vessel maintenance: manipulators, transfer casks, in-pipe
welding/cutting tools, and rad-hardened actuators. The building that
houses this equipment (hot cell) is costed separately in CAS21.

The cost model is:
\begin{equation}
\label{eq:cas2210}
C_{220110} = C_{\text{RH,base}}(f) \times S_{\text{concept}}
             \times \left(\frac{P_{\text{et}}}{1000}\right)^{0.5}
\end{equation}

\noindent where $C_{\text{RH,base}}(f)$ is a fuel-dependent base cost
(M\$ at 1~GWe tokamak reference) and $S_{\text{concept}}$ is a
geometry-dependent scaling factor.

\paragraph{Fuel dependence.}
Neutron activation determines whether remote handling is needed and how
heavily equipment must be rad-hardened.
\Cref{tab:rh-fuel} lists the base costs.

\begin{table}[htbp]
\caption{Remote handling base costs by fuel cycle (1~GWe tokamak reference).}
\label{tab:rh-fuel}
\centering
\begin{tabular}{lrl}
\toprule
Fuel & Base (M\$) & Rationale \\
\midrule
D-T          & 150 & Full RH suite, rad-hardened to 14.1~MeV neutrons \\
D-D          & 100 & Same scope, reduced rad-hardening (2.45~MeV) \\
D-${}^3$He   &  30 & Lighter duty, activation still exceeds occupational limits \\
p-${}^{11}$B &  20 & Conventional maintenance equipment (no rad-hardening) \\
\bottomrule
\end{tabular}
\end{table}

\paragraph{Concept dependence.}
Toroidal vessels (tokamak, stellarator) require articulated multi-axis
manipulators to reach components through narrow ports, while linear
geometries (mirror, FRC) permit simpler end-access tooling.
Toroidal concepts use $S_{\text{concept}} = 1.0$; linear concepts
use $S_{\text{concept}} = 0.55$.


%% ------------------------------------------------------------
\subsubsection{CAS22.01.12: Isotope Separation}
\label{sec:cas2212}

CAS22.01.12 is set to zero for all fuel cycles. All isotope procurement
is modeled as market purchase in CAS80 at enriched unit prices.

The rationale is that CAS22.01.12 (on-site separation plant capital) and
CAS80 (enriched market prices) are mutually exclusive models of the same
cost. If CAS80 pays \$2{,}175/kg for deuterium, that price already
includes the supplier's extraction and enrichment costs. Simultaneously
budgeting a \$300M on-site heavy-water plant charges for separation
twice. The correct model is one or the other:

\begin{itemize}
\item \textbf{Market purchase:} CAS80 at enriched \$/kg, CAS22.01.12 $= 0$.
\item \textbf{On-site plant:} CAS22.01.12 capital cost, CAS80 at raw
  feedstock prices.
\end{itemize}

We choose market purchase because on-site separation is not justified
at fusion-scale consumption rates. A 1~GWe D-T plant consumes
73--194~kg/yr of deuterium against a global enriched supply exceeding
80{,}000~kg/yr from heavy-water plants. Tritium is bred on-site in the
blanket and processed by the fuel handling system (CAS22.05); it is not a
market purchase. Lithium-6 enrichment capacity (${\sim}1$--$2$~t/yr
globally) could be stressed by a fleet of D-T reactors, but a
centralized enrichment facility shared across plants is a fleet-level
investment, not a per-plant CAS22 capital item.

\Cref{tab:fuel-prices} lists the retained CAS80 unit costs.

\begin{table}[htbp]
\caption{Fuel isotope unit costs (CAS80, market purchase).}
\label{tab:fuel-prices}
\centering
\begin{tabular}{lrl}
\toprule
Isotope & \$/kg & Basis \\
\midrule
Deuterium & 2{,}175 & STARFIRE (1980), inflation-adjusted \\
Li-6 (enriched) & 1{,}000 & 90\% enriched Li-6 \\
He-3 & 2{,}000{,}000 & Scarcity pricing \\
Protium & 5 & Commodity H$_2$ \\
B-11 (enriched) & 10{,}000 & FOAK estimate (B-10 tails) \\
\bottomrule
\end{tabular}
\end{table}


%% ------------------------------------------------------------
\subsubsection{CAS30: Capitalized Indirect Service Costs}
\label{sec:cas30}

CAS30 covers construction support services not attributable to specific
plant systems: field indirect costs (equipment rental, temporary
structures, consumables), construction supervision, and offsite design
services \citep{emwg2007}.

\paragraph{Literature review.}
Three methodologies exist in the fusion costing literature:

\begin{enumerate}
\item \textbf{Schulte (1978):} A fixed 35\% of total direct cost (TDC),
  split as 15\% construction supervision, 15\% design services,
  5\% field indirect costs \citep{schulte1978}.

\item \textbf{Miller/ARIES LSA fractions:} A Level of Safety Assurance
  (LSA) dependent fraction of TDC, ranging from 21.7\% at LSA=1
  (inherently safe, minimal nuclear-grade requirements) to 29.0\% at
  LSA=4 (full nuclear-grade construction) \citep{miller1998, piet1986,
  logan1985}. The LSA concept \citep{piet1986} captures the cost
  reduction from eliminating nuclear-grade construction requirements
  where the physics permits.

\item \textbf{PyFECONS bottom-up:} A power-scaling formula
  calibrated from a staffing model for a 150~MWe reference plant
  \citep{woodruff2026}. This yields 5.8\% of TDC at 1~GWe, reflecting
  the labor component; non-labor indirect costs (equipment rental,
  temporary structures, consumables) that dominate CAS31 scope are
  not included in the staffing model.
\end{enumerate}

No real fusion construction data exists to calibrate CAS30. Private
fusion companies do not publish cost breakdowns at this granularity;
ITER is a one-of-a-kind research facility; and UK~STEP and EU-DEMO
cost studies use ARIES-lineage models. For context, efficient
series-built fission plants (Korean APR1400, Chinese Hualong~One)
achieve indirect costs of 12--20\% of TDC, while troubled FOAK
projects (US Vogtle~3\&4) exceed 35\%.

\paragraph{Adopted model.}
We adopt a simpler aggregate model:

\begin{equation}
\label{eq:cas30}
\text{CAS30} = f_{\text{indirect}} \times \text{CAS20} \times
               \frac{T_c}{T_{\text{ref}}}
\end{equation}

\noindent where $f_{\text{indirect}} = \refFindirect$ (default), $T_c$ is the
construction time, and
$T_{\text{ref}} = \pgfmathprintnumber[precision=0]{\refTc}$~yr is a
reference duration. The construction-time scaling captures the primary driver
that varies between scenarios.

The default of 20\% is chosen as:
\begin{itemize}
\item Below Schulte's 35\% (1978 fission baseline, too high for NOAK fusion),
\item Below Miller's LSA=2 value of 23.2\% (ARIES tradition, likely
  somewhat high for NOAK),
\item Above PyFECONS's 5.8\% (labor-only staffing model),
\item Consistent with efficient NOAK nuclear construction internationally
  (15--20\% of TDC).
\end{itemize}

Sub-account precision (C31/C32/C35 splits) is not empirically
distinguishable for fusion and is not modeled.


%% ------------------------------------------------------------
\subsubsection{CAS60: Interest During Construction}
\label{sec:cas60}

Interest during construction (IDC) represents the cost of financing
the plant during the construction period, as introduced in
\cref{sec:capex-annual}.

\paragraph{Derivation.}
Assuming the total overnight cost $C$ is spent uniformly over $T_c$ years
with end-of-year disbursements at annual interest rate $i$, the future
value at commercial operation date is:

\begin{equation}
\text{FV} = \frac{C}{T_c} \sum_{k=0}^{T_c-1} (1+i)^k
           = \frac{C}{T_c} \cdot \frac{(1+i)^{T_c} - 1}{i}
\end{equation}

\noindent The IDC is the excess above the overnight cost:

\begin{equation}
\label{eq:idc}
\text{IDC} = C \left[ \frac{(1+i)^{T_c} - 1}{i\,T_c} - 1 \right]
           = f_{\text{IDC}}(i, T_c) \times C
\end{equation}

\noindent At reference conditions ($i = \refWACC$,
$T_c = \pgfmathprintnumber[precision=0]{\refTc}$~yr),
$f_{\text{IDC}} = \pgfmathprintnumber[precision=3]{\refIDCfrac}$.

\paragraph{Interaction with CAS90.}
CAS60 and CAS90 must use complementary conventions to avoid
double-counting construction-period financing. Two valid approaches
exist:

\begin{itemize}
\item \textbf{Option~A} (adopted): Explicit IDC in CAS60, plain CRF
  in CAS90.
  \begin{equation}
  \text{Total capital} = C + \text{IDC}, \quad
  \text{CAS90} = \text{CRF} \times \text{Total capital}
  \end{equation}

\item \textbf{Option~B}: No CAS60, effective CRF in CAS90.
  \begin{equation}
  \text{CAS90} = \text{CRF} \times (1+i)^T \times C
  \end{equation}
\end{itemize}

\noindent We adopt Option~A because it makes IDC visible as a separate
line item, uses the more realistic uniform-spending assumption (versus
lump-sum in Option~B), and is consistent with EMWG/ARIES cost reporting
conventions \citep{emwg2007}.

%% ------------------------------------------------------------
\subsubsection{CAS70: Annualized O\&M and Replacement}
\label{sec:cas70}

CAS70 comprises two sub-accounts:

\begin{itemize}
\item \textbf{CAS71} --- Annual operations and maintenance, scaled from
  a fuel-dependent base cost per MW of net electric capacity.
\item \textbf{CAS72} --- Annualized scheduled component replacement.
  Replacement costs are discounted to present value at the time of
  each replacement event and annualized via the CRF.
\end{itemize}

\paragraph{CAS71: Staffing-based O\&M.}
The CAS71 O\&M coefficient is derived from a bottom-up staffing model
that estimates headcount and compensation by division (operations,
maintenance, administration, technical, offsite) for each fuel type,
then adds non-labor costs (maintenance materials, insurance, supplies,
regulatory fees, waste treatment, and general overhead). The staffing
model is grounded in two reference datasets: the S-PRISM sodium fast
reactor staffing estimate \citep{boardman2000} (494~staff for a
1,520~MWe fission plant) and INL's first-principles SFR cost
estimation \citep{prosser2024}, which scales the S-PRISM baseline to
different plant sizes. Conventional gas CCGT staffing data (25--35
staff for $\sim$1~GWe) anchors the lower bound.

The key differentiator between fuels is neutron-related operational
overhead. All fusion machines are regulated under 10~CFR~Part~30
(byproduct material), not Part~50 (reactor licensing), eliminating
the armed security force, licensed operator requirements, emergency
planning zones, and nuclear QA programs that drive fission plant
staffing. However, fuels with higher neutron production require
radiation protection programs, radwaste management, and
radiation-controlled maintenance --- costs that scale roughly with
neutron flux. The resulting fuel-dependent coefficients are:

\begin{center}
\begin{tabular}{lrrrl}
\toprule
Fuel & Staff & O\&M (M\$/yr) & O\&M (k\$/MW/yr) & Key drivers \\
\midrule
DT   & 117 & 52.0 & 52 & Tritium systems, HP dept, radwaste, remote handling \\
DD   &  94 & 38.7 & 39 & Reduced neutron flux ($\sim\!\tfrac{1}{3}$ DT), smaller T inventory \\
DHe3 &  69 & 26.0 & 26 & $\sim$5\% neutron fraction, minimal tritium \\
pB11 &  59 & 23.6 & 24 & Aneutronic, no tritium, RSO only \\
\bottomrule
\end{tabular}
\end{center}

\noindent All values are at a 1~GWe reference in 2023~USD. For
comparison, S-PRISM fission O\&M is \$50k/MW/yr and modern gas CCGT
is $\sim$\$12k/MW/yr. The pB11 estimate (\$24k/MW/yr) sits at roughly
twice the conventional CCGT level, reflecting fusion-specific system
complexity (vacuum, magnets, plasma-facing components) rather than
regulatory burden.

\paragraph{Levelization.}
Both sub-accounts are levelized using the growing-annuity procedure
described in \cref{sec:lac}, which accounts for nominal cost escalation
over the full plant lifetime. A simpler approach that inflates costs
to first-year-of-operation dollars but neglects the continued
escalation over the remaining operating life would underestimate
levelized O\&M and fuel costs by approximately
\pgfmathprintnumber[precision=0]{\refUnderestimate}\% at reference
conditions ($g = \refG$,
$n = \pgfmathprintnumber[precision=0]{\refN}$~yr).


%% ------------------------------------------------------------
\subsubsection{CAS90: Annualized Financial Costs}
\label{sec:cas90}

CAS90 converts the total capital investment into an equivalent annual
payment:

\begin{equation}
\label{eq:cas90}
\text{CAS90} = \text{CRF}(i, n) \times \text{Total capital}
\end{equation}

\noindent where total capital includes IDC (CAS60). This uses the
plain CRF (\cref{eq:crf}), not the ``effective CRF''
$\text{CRF} \times (1+i)^{T_c}$ found in some references
\citep{miller1998}. As discussed in \cref{sec:cas60}, the effective
CRF is a valid convention only when CAS60 is zero; combining it with
an explicit CAS60 double-counts construction financing.

In a constant-dollar framework, CAS91 (operating-period escalation)
is zero by convention, and CAS92 (annual regulatory fees) is not
modeled. The framework thus equates CAS90 with CAS93 (cost of money).


%% ============================================================

\bibliographystyle{unsrtnat}
\bibliography{references}

\end{document}
